Dado el siguiente problema de programación lineal (P):


\begin{equation}
\begin{split}
\mbox{max. } z = 3x_1 + 2x_2\\
s.a.:\\
x_1 + 2x_2 \leq 60\\
x_1 + x_2 \geq 10\\
x_1 \leq 30\\
x_1,\,\,x_2\geq 0\\
\end{split}
%\label{ec:4}
\end{equation}

Se pide identificar una solución de cada uno de los siguientes tipos explicando por qué cumple con lo que se indica:
\begin{enumerate}
    \item Tipo 1: solución no básica y factible.
    \item Tipo 2: solución básica y factible.
    \item Tipo 3: solución básica y no factible.
    \item Tipo 4: solución no básica y no factible.
    \item ¿Cuál de los cuatro tipos de soluciones anteriores cumple que sus soluciones siempre ofrecen mejores valores de la función objetivo que las de cualesquiera otro de los tipos?
    \item Indica y justifica un valor que sea cota inferior del problema dual de (P).
\end{enumerate}